\subsection{A mixed cubic extension of the bar repair}

\begin{proposition}[C25: \status{Proved}; prescribed lower arities extend through arity three]
\label{sa08:C25}
Work over $\Z$, and let $P=K(\Z,2)$.
Let
\[
 U(P)=\Gamma(P[2])\otimes B\bigl(\Gamma(P[2])\bigr)
\]
be the positive bar closure of Statement~\ref{sa08:D029}.
An element with polynomial arity $m$ and bar length
$n$ has form degree $2m-n$. Give the original
parameter weight $w$.

Suppose a partial strict simplicial PD-DGA map
$U(P)\to A$ is already prescribed through
parameter arity two. Suppose that in weight $3w$
the target is vertically a resolution of
$\underline{\Z}$ through degree six, and its
rows zero through six are horizontally acyclic.
Then the prescribed partial map extends strictly
through arity three, without changing either
earlier arity.

More generally, for any simplicial $P$ with
degreewise-projective normalized terms, the
extension obstruction under these same target
bounds belongs to
\[
 \operatorname{Ext}^6_{s\mathrm{Ab}}
          (P/3,\underline{\Z}).
\]
For $P=K(\Z,2)$ this group is zero.
\end{proposition}
\begin{proof}
Let $V_3(P)\subset U_3(P)$ be the arity-three
part of the PD subalgebra generated by arities
one and two, and put $Q_3=U_3/V_3$.
The signed Lyndon freeness in
Statement~\ref{sa08:C24} implies that the
prescribed partial map induces a well-defined
map $V_3(P)\to A^{3w,*}$. To see why no
functorial basis is needed, choose the free
PD/exterior basis in each simplex separately.
The subalgebra generated by the generators
of arities at most two is free on those
generators. Thus their prescribed images
determine every decomposable expression
uniquely. This rule is intrinsic to the
lower-arity map, hence commutes with every
simplicial operator.

We compute the natural quotient complex.
All positive coefficient factors are
decomposable. In bar length two, the
shuffle relation for a weight-one and a
weight-two letter identifies their two
orders. In length three, products of a
single letter with a two-letter word
give the signed relations
\[
 abc-bac+bca.
\]
Consequently $Q_3$ is the complex
\begin{equation}
 \frac{P^{\otimes3}}
 {\langle abc-bac+bca\rangle}
 \ \xrightarrow{\delta}\
 P\otimes\Gamma^2P
 \ \xrightarrow{\mu}\
 \Gamma^3P
 \tag{B-C4}
\end{equation}
in form degrees $3,4,5$, where
\[
 \delta[abc]=a\otimes bc-c\otimes ab,
 \qquad \mu(a\otimes v)=av.
\]
A harmless sign change in the last term
puts the bar differential in this form.
The form-six coefficient term
$\Gamma^3P$ of $U_3$ is already in $V_3$;
it is not the remaining bar-letter
$\Gamma^3P$ in form five.

For a free abelian module choose a basis
and decompose (B-C4) by basis multidegree.
There are exactly three patterns.
\[
\begin{array}{c|c|c|c}
\text{pattern}&Q_3^3&Q_3^4&Q_3^5\\ \hline
aaa&0&\Z&\Z\\
aab&\Z&\Z^2&\Z\\
abc&\Z^2&\Z^3&\Z
\end{array}
\]
In the first row the nonzero differential
is multiplication by $3$.
In the second row the differentials are
\[
 \binom{1}{-2},\qquad (2\ \ 1).
\]
Indeed the word classes satisfy
$[aba]=0$ and $[baa]=-[aab]$, and
the form-four basis is
$a\otimes ab,\ b\otimes\gamma_2(a)$.
In the third row use $[abc],[acb]$
and the basis
$a\otimes bc,b\otimes ac,c\otimes ab$.
The differentials are
\[
 \begin{pmatrix}1&1\\0&-1\\-1&0\end{pmatrix},
 \qquad (1\ \ 1\ \ 1).
\]
The remaining word classes are
$[bac]=[abc]-[acb]$,
$[bca]=-[acb]$,
$[cab]=[acb]-[abc]$, and
$[cba]=-[abc]$.
The displayed columns form a primitive
basis of the kernel in each of the
last two rows: they have a unit maximal
minor. Thus the only vertical homology
is one $\Z/3$ in form five for each
diagonal pattern.

This is a natural identification
\begin{equation}
 Q_3(P)\simeq(P/3)[-5].
 \tag{B-C5}
\end{equation}
The quotient map is
$\gamma_3(a)\mapsto a\bmod3$, with
mixed products sent to zero.
The divided-power addition formula
and $n^3\equiv n\pmod3$ show that it
is additive, $\Z$-linear, and natural.
All normalized terms of $Q_3$ are
projective. The multidegree calculation
gives free simplicial terms, and
normalization is a direct summand.
Equivalently, the free construction
preserves retracts for projective
parameters.

Consider the exact sequence of cochain
complexes in $s\mathrm{Ab}$
\[
 0\longrightarrow V_3\longrightarrow U_3
       \longrightarrow Q_3\longrightarrow0.
\]
The obstruction to extending the given
chain map on $V_3$ lies in
$\operatorname{Ext}^1_{\operatorname{Ch}(s\mathrm{Ab})}
(Q_3,A^{3w,*})$.
Every extension splits in each cochain
degree, since the projective-normalized
terms of $Q_3$ have vanishing
$\operatorname{Ext}^1$ against the
horizontally acyclic target rows, as
proved before Statement~\ref{sa08:C22}.
A choice of these splittings identifies
the extension class with
\[
 H^1\operatorname{Tot}
       \operatorname{Hom}_{s\mathrm{Ab}}(Q_3,A^{3w,*}).
\]
Apply (B-C2) and (B-C3) to $Q_3[-1]$,
which is supported in form degrees
four through six. The stated finite
bounds give
\[
 H^1\operatorname{Tot}
       \operatorname{Hom}_{s\mathrm{Ab}}(Q_3,A^{3w,*})
 \cong
 \operatorname{Ext}^6_{s\mathrm{Ab}}
       (P/3,\underline{\Z}).
\]

For $P=K(\Z,2)$, normalization of
$P/3$ is $\Z/3$ in degree two.
Its projective dimension in the
abelian category of nonnegative
complexes is at most three, using
the identity-disk resolution and the
length-one free resolution of $\Z/3$.
The displayed Ext group is zero.
Therefore an actual simplicial
cochain map $U_3\to A^{3w,*}$
extending the prescribed decomposable
map exists.

Every multiplication or divided-power
identity whose output has arity three
already belongs to $V_3$. Hence this
chain extension, together with the
unchanged lower arities, is precisely
a partial strict simplicial PD-DGA
map through arity three.
\end{proof}

This result permits arbitrary earlier choices
through arity two. It therefore treats mixed
terms that the single-disk factorization in
Statement~\ref{sa08:C23} does not cover.
It does not provide all higher arities,
prescribed parameter-boundary extensions,
or a uniform weight satisfying infinitely
many eventual bounds.

The exact certificate
\texttt{B\_bar\_arity3\_checks.py} records
the three multidegree calculations in
\texttt{B\_bar\_arity3\_checks.json}.
It verifies all signed shuffle relations,
explicit quotient sections, unimodular
relation bases, the induced bar
differentials, and primitive kernel
images. These three patterns exhaust
the free-module calculation, rather
than sampling dimensions.
